\chapter{Por qué el espacio de secuencias}
\label{cap:porque}

\section{El problema: la distribución como variable de estado}

Empecemos por entender exactamente qué se rompe cuando uno agrega
heterogeneidad a un modelo dinámico.

En un modelo de agente representativo, el estado del sistema en la fecha $t$ es
un vector de pocos números: el capital, quizá la deuda pública, el nivel de
productividad. Llamémoslo $x_t \in \R^n$ con $n$ pequeño, digamos menor que
treinta. Resolver el modelo a primer orden consiste en encontrar matrices
$A$ y $B$ tales que
\begin{equation}
x_t = A\, x_{t-1} + B\, \epsilon_t ,
\end{equation}
y eso lo hace Dynare en milisegundos desde hace veinte años.

Ahora pongamos un continuo de hogares que difieren en su riqueza $a$ y en su
productividad $e$, que no pueden endeudarse y que enfrentan riesgo idiosincrático
no asegurable. ¿Cuál es el estado del sistema? Ya no basta con el capital
agregado: para saber cuánto ahorrará mañana la economía hay que saber
\emph{cómo está repartida} la riqueza hoy, porque un hogar pobre y uno rico
responden distinto al mismo cambio en la tasa de interés. El estado es ahora la
\emph{distribución completa} $D_t(e,a)$, que es un objeto de dimensión infinita.
Al discretizarla sobre una grilla de, digamos, 7 estados de productividad y 200
puntos de activos, el estado tiene $7 \times 200 = 1\,400$ dimensiones. En
modelos de dos activos, con dos estados endógenos, es fácil llegar a $10\,000$ o
$250\,000$.

Ese es el problema, y tiene un nombre: la \emph{maldición de la
dimensionalidad} del espacio de estados.

\section{Tres respuestas históricas}

\paragraph{Krusell y Smith (1998)} La primera respuesta famosa fue conjeturar
que, aunque la distribución sea de dimensión infinita, los agentes solo necesitan
seguir unos pocos de sus momentos --- típicamente la media --- para pronosticar
los precios con suficiente precisión. Se resuelve el modelo suponiendo esa
``racionalidad acotada'', se simula, se verifica que el $R^2$ de la ley de
movimiento conjeturada sea altísimo y se itera. Funciona notablemente bien en el
modelo original, pero es lento, no da una solución exacta del modelo linealizado
y es difícil de usar para estimación: hay que repetir toda la simulación para
cada vector de parámetros.

\paragraph{Reiter (2009)} La segunda respuesta fue: linealicemos solo en los
agregados. Los agentes resuelven su problema individual de forma global y exacta,
pero las desviaciones agregadas respecto del estado estacionario se tratan a
primer orden. Formalmente se escribe un sistema lineal en el espacio de estados,
donde el estado incluye la distribución discretizada completa. Esto es exacto
(a primer orden en agregados) y rápido de evaluar una vez construido, pero
construirlo requiere trabajar con matrices de tamaño $n_g \times n_g$ donde
$n_g$ es el número de puntos de la grilla, y el costo escala con el \emph{cubo}
de ese número. Por eso la literatura posterior dedicó mucho esfuerzo a
\emph{reducir} el modelo antes de resolverlo (Ahn \emph{et al.}, 2018;
Winberry, 2018; Bayer y Luetticke, 2020).

\paragraph{Boppart, Krusell y Mitman (2018)} La tercera respuesta cambió el
espacio. En lugar de preguntarse ``¿cuál es la ley de movimiento del estado?'',
se preguntaron ``¿cuál es la respuesta del modelo a un choque inesperado que
perturba el estado estacionario?''. Esa respuesta --- la función impulso-respuesta
de previsión perfecta, o choque \emph{MIT} --- se puede calcular resolviendo una
transición determinista, sin tocar nunca la ley de movimiento del estado. Y por
equivalencia cierta, a primer orden esa función impulso-respuesta \emph{es} la
representación de medias móviles del modelo con riesgo agregado.

\section{La idea de \ABRS}

\ABRS\ parten de esta tercera respuesta y la llevan hasta sus últimas
consecuencias. Si lo que uno quiere son funciones impulso-respuesta, y si el
modelo es lineal en agregados, entonces todo el modelo queda descrito por un
conjunto de \emph{derivadas de secuencias respecto de secuencias}.

Pongámoslo concreto. En un modelo de mercados incompletos, el ahorro agregado de
los hogares en la fecha $t$ depende de toda la trayectoria presente y futura de
tasas de interés y salarios:
\begin{equation}
A_t = \mathcal{A}_t\big(\{r_s, w_s\}_{s\ge 0}\big).
\end{equation}
Esta función esconde todo: cómo responde cada hogar, cómo evoluciona la
distribución, cómo interactúan. Pero si solo nos interesa el comportamiento
\emph{local} alrededor del estado estacionario, lo único que necesitamos de esa
función es su derivada, es decir, las matrices
\begin{equation}
\J^{A,r}_{t,s} \equiv \frac{\partial A_t}{\partial r_s}\bigg|_{\sst},
\qquad
\J^{A,w}_{t,s} \equiv \frac{\partial A_t}{\partial w_s}\bigg|_{\sst}.
\end{equation}

\begin{idea}
$\J^{A,r}$ es un \emph{estadístico suficiente} del bloque heterogéneo para el
equilibrio general. Dos modelos con distribuciones de riqueza completamente
distintas pero con el mismo $\J^{A,r}$ producen exactamente las mismas funciones
impulso-respuesta agregadas. Toda la complejidad de la heterogeneidad se
condensa en esa matriz.
\end{idea}

Esto tiene tres consecuencias prácticas enormes.

\begin{enumerate}[leftmargin=1.4em]
\item \textbf{El tamaño del sistema deja de depender del espacio de estados.}
Las matrices son $T\times T$ con $T\approx 300$, independientemente de si la
grilla tiene 1\,400 o 250\,000 puntos. En la tabla 1 de \abrs\ se ve que el
tiempo de resolución del equilibrio general es prácticamente idéntico entre una
calibración con $3\,500$ puntos y otra con $250\,000$.

\item \textbf{Los jacobianos se reciclan.} Al estimar un modelo hay que evaluar
la verosimilitud miles de veces. Si los parámetros que uno estima no cambian el
estado estacionario del problema individual (persistencia de los choques,
rigidez de precios, coeficientes de la regla de política), entonces el jacobiano
del bloque heterogéneo \emph{no cambia}: se calcula una vez y se reutiliza. Lo
que queda por hacer en cada evaluación es álgebra lineal de milisegundos.

\item \textbf{El modelo se vuelve modular.} Cada pieza de la economía --- las
empresas, el gobierno, los bancos, los hogares --- es un bloque con insumos y
productos. Los bloques se componen con la regla de la cadena a lo largo de un
grafo. Cambiar la regla de política monetaria no obliga a recalcular nada del
bloque de hogares.
\end{enumerate}

\section{Lo que se gana y lo que se pierde}
\label{sec:limites-perturbacion}

Conviene ser explícito sobre el precio. El método es una \emph{perturbación de
primer orden en los agregados}. Eso implica, igual que para Reiter y para
cualquier método de esta familia:

\begin{itemize}[leftmargin=1.4em]
\item \textbf{No hay primas de riesgo.} Por equivalencia cierta, los agentes se
comportan como si no hubiera incertidumbre agregada. El riesgo
\emph{idiosincrático} sí está presente de forma completamente no lineal --- es
lo que genera la distribución --- pero el riesgo \emph{agregado} desaparece a
primer orden.
\item \textbf{La elección de cartera agregada es indeterminada.} Si dos activos
tienen el mismo retorno a primer orden, el modelo no dice cuánto se tiene de
cada uno.
\item \textbf{La política óptima de Ramsey está mal definida}, porque requiere
evaluar el bienestar a segundo orden.
\item \textbf{Los efectos no lineales de choques grandes se pierden}, aunque
el capítulo \ref{cap:nolineal} muestra cómo recuperarlos con un método de
cuasi-Newton que reutiliza los mismos jacobianos.
\end{itemize}

\begin{tutesis}
Esta limitación importa para tu pregunta. Un modelo de seguro de liquidez en el
que el ``estado agregado de escasez de dólares'' es un evento de cola --- poco
probable y muy costoso --- es justamente un modelo donde las primas de riesgo
podrían ser el objeto de interés. Hay una salida, y la usamos en el capítulo
\ref{cap:aplicacion}: mantener la incertidumbre agregada \emph{dentro} del
período, como un conjunto finito de estados con probabilidades fijas que se
resuelven simultáneamente, y perturbar solo la secuencia de esas probabilidades
y magnitudes. Así el modelo sigue siendo lineal en la secuencia pero conserva
el no linealidad del evento de cola. Es exactamente la estructura de tus
ecuaciones (11)--(19).
\end{tutesis}

\section{Mapa del método en una página}

El procedimiento completo, que el resto del manual desarrolla, tiene cinco pasos.

\begin{receta}
\begin{enumerate}[leftmargin=1.4em,itemsep=2pt]
\item \textbf{Resolver el estado estacionario} del problema individual y
encontrar la distribución estacionaria. Esto es un problema de punto fijo
estándar, no tiene nada de nuevo, y se cubre en el capítulo \ref{cap:bloque}.
\item \textbf{Calcular los jacobianos del bloque heterogéneo} con el algoritmo
fake news: una sola iteración hacia atrás, unos pocos productos de matrices y una
recursión acumulativa. Capítulo \ref{cap:fakenews}.
\item \textbf{Armar el modelo como un grafo} de bloques simples y bloques
heterogéneos, declarar incógnitas y objetivos, y componer los jacobianos por la
regla de la cadena. Capítulo \ref{cap:dag}.
\item \textbf{Resolver el equilibrio general}: $d\mathbf{U} = -\mathbf{H}_U^{-1}
\mathbf{H}_Z\, d\mathbf{Z}$. Una inversión de matriz. Capítulo \ref{cap:dag}.
\item \textbf{Usar la solución}: funciones impulso-respuesta, descomposiciones,
momentos, verosimilitud, estimación bayesiana, transiciones no lineales.
Capítulos \ref{cap:usos} a \ref{cap:nolineal}.
\end{enumerate}
\end{receta}

Los pasos 1 y 2 son los únicos que cuestan tiempo de cómputo. Los pasos 3 a 5
son milisegundos. Esa asimetría es la que hace viable la estimación.
